\section{Related Work}\label{sec:related}

\paragraph{Nonstationary optimization and dynamic regret.}
Dynamic-regret and variation-budget analyses show that under drift, forgetting
scales and window lengths must adapt to how quickly the target moves
\citep{BesbesGurZeevi2015,CheungSimchiLeviZhu2019,YuanLamperski2020,
MazzettoUpfal2023,WangXieChenLuiZhou2024}. That literature isolates the control
importance of retention scale. The horizon law gives the corresponding
statistical law for distribution tracking: finite memory has a temporal-validity
horizon determined jointly by a finite-sample exponent and a drift roughness
exponent.

This distinction is especially sharp relative to adaptive update policies such as
\citet{MazzettoUpfal2023}. Those works optimize a decision or intervention rule
under evolving data. The present object is earlier in the pipeline: it asks how
long retained evidence continues to represent the present distribution before age
becomes the dominant source of tracking error.

\paragraph{Adaptive windowing, drift detection, and update policies.}
Adaptive-window and concept-drift methods study when accumulated evidence is
sufficient to react \citep{BifetGavalda2007,Gama2004,GamaSurvey2014}.
\citet{hinder2023window} emphasize that the window length itself is part of the
problem. Resource-aware update scheduling under drift addresses a related
decision problem in which update cost must be balanced against accumulated loss
\citep{PiasecznyEtAl2025RCCDA}, and finite-time analyses of sliding-window or
discounted adaptation under switching environments make the same design variable
explicit from the regret side \citep{GarivierMoulines2011}. Streaming OT and MMD detectors contribute the
alarm-time perspective, emphasizing delay, power, and memory usage
\citep{ShvetsovEtAl2020OTChangePoint,KalinkeEtAl2025MMDEW}. These works provide
the reaction timescale. The horizon law adds the validity timescale: retained
evidence can cease to represent the present before change statistics cross an
alarm threshold.

\paragraph{Adaptive filtering and locally stationary smoothing.}
Classical adaptive filtering and state-space tracking tune forgetting factors,
gain schedules, or process-noise levels to balance smoothing against
responsiveness \citep{Kalman1960,Jazwinski1970}. Locally stationary smoothing
and bandwidth selection solve the analogous bias--variance problem for evolving
regression or spectral targets \citep{Dahlhaus1997,Vogt2012}. Recent Wasserstein-
based results for locally stationary functional time series move that same local
windowing logic directly into a transport metric \citep{TinioAlayaBouzebda2025}.
These literatures supply mechanisms for adaptation. The horizon law provides the distributional law
that such mechanisms must respect when the retained object is an empirical
distribution rather than a latent state or a local regression surface.

Lag-geometry inference under noisy transport observations is complementary to
that structural law. 
\citet{Parreira2026Jul} study power-law lag geometry through a parameter-coupled
heteroscedastic regression, deriving an information law for estimating
$(H,\zeta)$ and testing the adequacy of the lag model itself. That inferential
layer supplies one route for estimating the roughness inputs that enter the
horizon law developed here.

\paragraph{Distribution tracking under transport geometry.}
Wasserstein distance and Sinkhorn-type estimators provide the finite-sample
geometries that enter the horizon law. For empirical Wasserstein distances,
\citet{FournierGuillin2015}, \citet{BoissardLeGouic2014}, and
\citet{NilesWeedBach2019} show that the sampling exponent depends on geometric
complexity even before drift enters. The horizon law takes that finite-sample
exponent as one input and composes it with a staleness exponent induced by the
drift path. In this sense, the transport literature provides the sampling side
of the law, while the horizon law provides its temporal counterpart.

\paragraph{Entropic and smoothed transport as operational geometries.}
Regularized transport changes both statistical stability and computational
behavior \citep{GenevayEtAl2019,MenaWeed2019,delBarrioEtAl2023,
RigolletStromme2025,Stromme2023MID,GroppeHundrieser2024LCA,
EcksteinNutz2022,GhosalNutzBernton2022}. Fixed-$\varepsilon$ Sinkhorn theory
provides i.i.d. asymptotics, differentiability, null limit laws, and stability
constants \citep{GoldfeldEtAl2022}. Those results are the finite-sample anchor
for the operational horizon law developed here. Recent linear-Gaussian entropic-
OT analyses make the recursive structure of regularized transport explicit
\citep{AkyildizEtAl2024}, suggesting a constructive intermediate regime between
fixed-$\varepsilon$ i.i.d. asymptotics and fully drifted triangular-array
tracking.

Smooth-functional and delta-method arguments provide the multivariate transfer
mechanism behind this regime. Hadamard differentiability is the natural regularity
class for lifting triangular-array fluctuations through a transport functional,
and the fixed-$\varepsilon$ Sinkhorn map sits in that category after gauge
fixing and uniform stability control \citep{GoldfeldEtAl2022}. That route is the
right multivariate complement to the one-dimensional quantile-process proof
model.

\paragraph{Intrinsic geometry and operational calibration.}
Geometry-aware calibration of regularization scales points toward intrinsic
dimension as the relevant complexity variable in regularized transport
\citep{GenansWintenberger2026}. Local intrinsic-dimension diagnostics such as
TwoNN provide a direct statistic for effective support geometry
\citep{FaccoEtAl2017}. That perspective enters through the operational Sinkhorn
frontier: intrinsic geometry helps identify which finite-sample regime is
available, and therefore which horizon law is relevant.

\paragraph{Low-dimensional surrogate geometries.}
Sliced, max-sliced, subspace-robust, and projection-robust Wasserstein
constructions replace ambient transport by one-dimensional or low-dimensional
comparisons \citep{KolouriRohdeHoffmann2018,DeshpandeEtAl2019,PatyCuturi2019,
LinEtAl2020,NadjahiEtAl2020,NietertEtAl2022}. These examples show concretely
that geometry choice can change the finite-sample exponent. They therefore fit
naturally into the horizon framework as alternative sampling geometries whose
drift-induced useful-memory scales need not match raw Wasserstein behavior.

\paragraph{Wasserstein nonstationarity and Age of Information.}
Wasserstein nonstationarity measures already appear in online optimization and
distributionally robust optimization under drift
\citep{JiangLiZhang2020,KeehanAndersonWiesemann2025,AminiVinas2026}. The
Age-of-Information literature treats age as a resource that must be balanced
against communication or update costs \citep{KaulYatesGruteser2012,
YatesEtAl2021AoI}. Together these viewpoints make age and nonstationarity into
operational variables. The horizon law gives them a distributional statistical
interpretation: age is costly because it accumulates geometric mismatch with the
present, and the roughness exponent determines how fast that cost grows.

\citet{KeehanAndersonWiesemann2025} are particularly close in spirit: they
derive concentration bounds for weighted empirical laws under sequential
Wasserstein drift and optimize weights for distributionally robust decisions.
That question begins once a decision rule is fixed. The temporal-validity
horizon addresses the preceding statistical question of how long evidence should
be retained before staleness dominates the finite-sample gain. In this sense,
their weighting problem and the present horizon problem are complementary rather
than competing.

Weighted empirical-process and weighted-Wasserstein precursors are also relevant
for the non-uniform-memory line. In particular, \citet{Mason2016} studies
weighted approximations to empirical Wasserstein fluctuations, providing a
technical ancestor for the weighted-memory summaries $n_{\mathrm{eff}}(w)$ and
$M_H(w)$ used here.

Across these neighboring literatures, the shared theme is that
nonstationarity makes retention scale matter. The horizon law identifies the
statistical object that organizes that dependence for distribution tracking: a
temporal-validity horizon, together with a useful-memory region and a distinct
detectability timescale.
